
\clearpage
\pageoriginale

\SGAExposeBookmark{XIII}{El formalismo de los ciclos evanescentes}
\begin{center}
{\large\textbf{SGA 7}}\\[1.5em]
{\Large\underline{EXPOSÉ XIII}}\\[2.5em]
{\Large\underline{EL FORMALISMO DE LOS CICLOS EVANESCENTES}}\\[1.5em]
{\large\underline{por P. DELIGNE}}
\end{center}

\bigskip
\noindent\textbf{Índice de la fuente}
\begin{center}
\begin{tabular}{@{}llr@{}}
Introducción && 2\\
\S1. & Haces de conjuntos & 6\\
1.1 & Haces galoisianos & 6\\
1.2 & El topos $Y\times_s S$ & 8\\
1.3 & El funtor $\Psi$ & 13\\
1.4 & La variación & 16\\
\S2. & Ciclos evanescentes & 17\\
2.1 & El funtor $R\Psi$ & 17\\
2.2 & Una compatibilidad (morfismos traza) & 26\\
2.3 & Teorema de finitud (en característica igual a cero) & 29\\
2.4 & Singularidades aisladas & 30
\end{tabular}
\end{center}

\pageoriginale

\bigskip
\noindent\underline{Introducción}

En esta exposición retomamos el formalismo introducido en el n.º~2 del Exposé~I
de Grothendieck. El único punto nuevo es la introducción del análogo
algebraico de la \underline{variación} de la teoría clásica.

La imagen topológica que tengo presente es la siguiente. Sea $X$ un espacio
analítico, sea $D\subset\mathbb C$ un disco de centro $0$, pongamos
$D^*=D\setminus\{0\}$ y sea $f:X\longrightarrow D$ un morfismo propio (el
caso no propio se considerará más adelante). Reemplacemos $D$ por un disco
suficientemente pequeño de centro $0$, reemplacemos $X$ por la imagen inversa
de ese pequeño disco y sigamos denotando el resultado por
$f:X\longrightarrow D$. Según la variante analítica compleja del teorema de
isotopía de Thom, $f^{-1}(D^*)$ es un fibrado topológico sobre $D^*$. No sé
si siempre puede construirse un par $(\Gamma,r)$ del tipo que se describe a
continuación, pero parece verosímil que así sea.

\begin{enumerate}[label=(\alph*),leftmargin=2.6em]
\item $\Gamma$ es un sistema localmente coherente de trivializaciones locales
del fibrado $f^{-1}(D^*)$ sobre $D^*$. En otras palabras, $\Gamma$ define una
trivialización $\pi_1(D^*)$-equivariante de la imagen inversa de $X$ sobre un
recubrimiento universal $\widetilde D^*$ de $D^*$.

\item $r$ es una retracción de $X$ sobre la fibra especial
$X_0=f^{-1}(0)$, compatible con $\Gamma$: si, sobre un abierto conexo $U$ de
$D^*$, $\Gamma$ identifica $X_U=f^{-1}(U)$ con $X_1\times U$, entonces
$r=r(x_1,u)$ es independiente de $u$.
\end{enumerate}

Supongamos dado $(X,f,\Gamma,r)$. Sea $t\in D$ y sea
$X_t=f^{-1}(t)$ la «fibra general». El fibrado $\varphi^*X$, obtenido de
$X/D$ por cambio de base mediante
\begin{equation*}
\varphi:[0,1]\longrightarrow D^*,
\qquad x\longmapsto \exp(2\pi i x)\,t,
\end{equation*}
queda trivializado por $\Gamma$. La \underline{monodromía}
$T:X_t\longrightarrow X_t$ es el homeomorfismo
\begin{equation*}
X_t=(\varphi^*X)_0\xrightarrow[\Gamma]{\sim}
(\varphi^*X)_1=X_t.
\end{equation*}
Sea $r_t:X_t\longrightarrow X_0$ la restricción de $r$ a $X_t$. Entonces
$r_tT=r_t$. Se reconstruye $(X,f)$ a partir de $(X_t,X_0,T,r_t)$ como sigue:
\begin{enumerate}[label=(\alph*),leftmargin=2.6em]
\item se identifican mediante $T$ los extremos $X_t\times\{0\}$ y
$X_t\times\{1\}$ de $X_t\times[0,1]$, obteniéndose
$f':X'\longrightarrow S^1$;
\item $r_t$ define $r':X'\longrightarrow X_0$, y $(X,f)$ se identifica con
el morfismo de conos
\begin{equation*}
\operatorname{cone}(f'):
\operatorname{cone}(r':X'\longrightarrow X_0)
\longrightarrow D=\operatorname{cone}(S^1\longrightarrow\mathrm{pt}).
\end{equation*}
\end{enumerate}

\pageoriginale

Pongamos $\Psi=(r,f):X\longrightarrow X_0\times D$. Las imágenes directas
superiores de Leray $R^i\Psi_*\mathbb Z$ satisfacen la condición siguiente:

\begin{quote}
(*) El espacio $D^*$ está recubierto por abiertos $U$ tales que la
restricción del haz $F$ a $X_0\times U$ proviene, por imagen inversa, de un
haz sobre $X_0$.
\end{quote}

\medskip
\noindent\underline{Construcción 0.1.}\quad
\underline{Los haces $F$ sobre $X_0\times D$ que satisfacen (*) se
identifican con los sistemas $(F_0,F_t,T,\alpha)$ en los que}
\begin{enumerate}[label=\alph*),leftmargin=2.5em]
\item \underline{$F_0$ es un haz sobre $X_0$};
\item \underline{$F_t$ es un haz sobre $X_0$, provisto de una acción del
generador $T$ de $\pi_1(D^*,t)$};
\item \underline{$\alpha$ es un morfismo de $F_0$ al haz de invariantes por
$T$ de $F_t$}.
\end{enumerate}

Sea $G$ un haz sobre $X_0\times D^*$ que satisface (*), y sea
$\widetilde D^*$ el recubrimiento universal de $D^*$, con punto base $t$.
La imagen inversa de $G$ sobre $X_0\times\widetilde D^*$ proviene, de una y
solo una manera, de un haz $G_t$ sobre $X_0$. El haz $G_t$ está provisto de
una acción de $\pi_1(D^*,t)$, y $G$ se obtiene de
$\operatorname{pr}_1^*G_t$ sobre $X_0\times\widetilde D^*$ tomando el
cociente por $\pi_1(D^*,t)$.

Un haz $F$ sobre $X_0\times D$ queda determinado por
\begin{enumerate}[label=\arabic*),leftmargin=2.5em]
\item su restricción $F_0$ a $X_0$;
\item su restricción $G$ a $X_0\times D^*$;
\item un morfismo $\alpha:F_0\longrightarrow i^*j_*G$, donde
$i:X_0\hookrightarrow X_0\times D$ y
$j:X_0\times D^*\hookrightarrow X_0\times D$.
\end{enumerate}
Si $F$ satisface (*), entonces $(G_t,T)$ describe $G$, y de ello se sigue la
descripción afirmada de $F$.

Para los haces $R^i\Psi_*\mathbb Z$ se tiene
\begin{equation*}
(R^i\Psi_*\mathbb Z)_0=
\begin{cases}
\mathbb Z, & i=0,\\
0, & i>0,
\end{cases}
\end{equation*}
mientras que, para $i>0$, los $(R^i\Psi_*\mathbb Z)_t$ merecen el nombre de
\underline{haces de ciclos evanescentes} (véase la discusión al comienzo del
párrafo~2 del Exposé~I).

\pageoriginale

Esta discusión solo resulta útil después de afinarla pasando a la categoría
derivada: los $R^i\Psi_*\mathbb Z$ se obtienen de un objeto de
$D^+(X_0\times D)$ y, más precisamente, de un objeto
$R\Psi_*\mathbb Z$ de la categoría derivada de la categoría de los haces
abelianos sobre $X_0\times D$ que satisfacen (*). Este objeto determina no
solo los haces de ciclos evanescentes y su monodromía, sino también los
morfismos clásicos de «variación» (véanse 1.4 y XIV, 3.1).

Definiremos su análogo en geometría algebraica. Si $S$ es un rasgo
estrictamente henseliano y $f:X\longrightarrow S$ es un morfismo de esquemas,
definiremos un objeto $R\Psi(\mathbb Z/\ell^n)$ de la categoría derivada de
la categoría de los haces abelianos sobre el topos producto
$(X_0)_{\mathrm{et}}\times S_{\mathrm{et}}$. Un haz sobre este topos puede
interpretarse como un sistema $(F_0,F_t,T,\alpha)$ en el que
\begin{enumerate}[label=\alph*),leftmargin=2.5em]
\item[$\alpha$)] $F_0$ es un haz sobre $X_0$;
\item[$\beta$)] $F_t$ es un haz sobre $X_0$, provisto de una acción continua $T$ del
grupo de inercia $I$;
\item[$\gamma$)] $\alpha$ es un morfismo de $F_0$ al haz de invariantes por $I$ de
$F_t$.
\end{enumerate}

Tendremos que trabajar en un marco un poco más general que el descrito arriba:
\begin{enumerate}[label=\alph*),leftmargin=2.5em]
\item La construcción de $R\Psi(\mathbb Z/\ell^n)$ no exige que $f$ sea
propio. Para una interpretación topológica se repetiría la construcción
precedente tras reemplazar $X$ por un entorno de la fibra especial $X_0$.

\item Para construir $R\Psi(\mathbb Z/\ell^n)$ y demostrar sus propiedades,
hay que considerar no solo el haz constante $\mathbb Z/\ell^n$, sino también
haces generales e incluso complejos de haces.

\item Cuando no se supone separablemente cerrado el cuerpo residual $k$ del
rasgo henseliano $S$, hay que reemplazar el topos producto
$(X_0)_{\mathrm{et}}\times S_{\mathrm{et}}$ por el topos producto fibrado
\begin{equation*}
(X_0)_{\mathrm{et}}
\times_{\operatorname{Spec}(k)_{\mathrm{et}}}S_{\mathrm{et}}.
\end{equation*}
El lector al que desagraden los productos $2$-fibrados de topoi puede tomar
como definición de este topos la descripción galoisiana de 1.2.4.
\end{enumerate}

\pageoriginale

\medskip
\noindent 0.2. \underline{Notación y terminología}.

\begin{description}
\item[(0.2.1) \underline{rasgo}:]
el espectro de un anillo de valoración discreta.

\item[(0.2.2) \underline{henseliano}:]
para un esquema, el espectro de un anillo local henseliano.

\item[(0.2.3) \underline{estrictamente henseliano}:]
para un anillo, henseliano con cuerpo residual separablemente cerrado; para
un esquema, el espectro de tal anillo.

\item[(0.2.4) \underline{punto geométrico} de $S$:]
un morfismo $\bar x:\operatorname{Spec}(\bar k)\longrightarrow S$ cuya imagen
es $x\in S$, donde $k(\bar x)=_{\mathrm{def}}\bar k$ es una clausura
separable de $\kappa(x)$. Se dice que $\bar x$ está localizado en $x$. Por
abuso de lenguaje, llamaremos a veces también punto geométrico a un morfismo
$\bar x:\operatorname{Spec}(\bar k)\longrightarrow S$ cuando $\bar k$ sea
una extensión separablemente cerrada de $\kappa(x)$; tal morfismo se
factoriza de manera única a través de un punto geométrico de $S$ en sentido
propio.

\item[(0.2.5) $s,\eta,\bar s,\bar\eta,\bar S$:]
Si $S=\operatorname{Spec}(V)$ es un rasgo henseliano, denotaremos normalmente
por $s$ el punto cerrado de $S$, por $\eta$ su punto genérico, por $\bar\eta$
un punto geométrico genérico de $S$ (localizado en $\eta$), y por $\bar s$ el
punto geométrico correspondiente localizado en $s$ (I, 0.0.3). Sea $\bar S$
el espectro de la normalización de $V$ en $k(\bar\eta)$; su cuerpo residual es
una extensión inseparable de $k(\bar s)$. Escribiremos
$\operatorname{Gal}(\bar\eta/\eta)$ para el grupo de Galois de
$k(\bar\eta)$ sobre $k(\eta)$, y análogamente para $s$. Se tiene una sucesión
exacta
\begin{equation*}
0\longrightarrow I\longrightarrow
\operatorname{Gal}(\bar\eta/\eta)\longrightarrow
\operatorname{Gal}(\bar s/s)\longrightarrow0,
\end{equation*}
donde $I$ es el grupo de inercia.

\item[(0.2.6) \underline{haz}:]
sobre un esquema $X$, significa siempre un haz sobre el sitio étale
$X_{\mathrm{et}}$ de $X$.

\item[(0.2.7) \underline{primo a}:]
un haz de torsión $F$ sobre un esquema $S$ (o un grupo abeliano de torsión
$F$) es primo a las características residuales de $S$ si, para todo
$s\in S$, la multiplicación por el exponente característico de $k(s)$ es un
automorfismo de $F$.
\end{description}

\pageoriginale

\bigskip
\noindent 1. \underline{Haces de conjuntos}

Los haces de ciclos evanescentes se definirán como valores de funtores
derivados del funtor exacto a izquierda $\Psi$, definido en el n.º~3 para
todo haz de conjuntos.

\medskip
\noindent 1.1. \underline{Haces galoisianos}.

\medskip
\noindent 1.1.1. Sea $Y$ un esquema sobre un cuerpo $k$ de clausura separable
$\bar k$, y pongamos $\bar Y=Y\otimes_k\bar k$. Sea $G$ un grupo profinito y
sea
\begin{equation*}
u:G\longrightarrow\operatorname{Gal}(\bar k/k)
\end{equation*}
un homomorfismo continuo de $G$ al grupo de Galois de $\bar k$ sobre $k$. El
grupo de Galois $\operatorname{Gal}(\bar k/k)$ actúa por la izquierda sobre
$\bar Y$ por transporte de estructura. El grupo $G$ actúa sobre $\bar Y$ a
través de $u$.

Sea $\mathcal F$ un haz de conjuntos sobre (el sitio étale de) $\bar Y$. Una
acción de $G$ sobre $\mathcal F$ \underline{compatible} con la acción de $G$
sobre $\bar Y$ es una acción de $G$ por automorfismos del par
$(\bar Y,\mathcal F)$ que induce la acción dada sobre $\bar Y$. En otras
palabras, es un sistema de isomorfismos
\begin{equation*}
\sigma(g):u(g)_*\mathcal F\longrightarrow\mathcal F
\end{equation*}
que satisface $\sigma(gh)=\sigma(g)\,\sigma(h)$. Si $G$ actúa sobre
$\mathcal F$, entonces, para todo $U$ étale sobre $Y$, con
$\bar U=U\otimes_k\bar k$ étale sobre $\bar Y$, el grupo $G$ actúa sobre
$\mathcal F(\bar U)$.

\medskip
\noindent\underline{Definición 1.1.2.}\quad
\underline{Diremos que $G$ actúa continuamente sobre $\mathcal F$ si, para
todo $U$ cuasicompacto y étale sobre $Y$, la acción de $G$ sobre el conjunto
discreto $\mathcal F(\bar U)$ es continua.}

Sería equivalente considerar únicamente los $U$ afines.

Si $\mathcal F$ es un haz de conjuntos sobre $Y$, con imagen inversa
$\bar{\mathcal F}$ sobre $\bar Y$, entonces
$\operatorname{Gal}(\bar k/k)$ actúa sobre $\bar{\mathcal F}$ por transporte
de estructura, compatiblemente con su acción sobre $\bar Y$.

\medskip
\noindent\underline{Recordatorio 1.1.3.}
\begin{enumerate}[label=(\roman*),leftmargin=2.7em]
\item \underline{La acción de $\operatorname{Gal}(\bar k/k)$ sobre
$\bar{\mathcal F}$ es continua.}

\pageoriginale

\item \underline{El funtor}
\begin{equation*}
\mathcal F\longmapsto\bar{\mathcal F},
\end{equation*}
\underline{donde $\bar{\mathcal F}$ está provisto de la acción de
$\operatorname{Gal}(\bar k/k)$, es una equivalencia de la categoría de los
haces de conjuntos sobre $Y$ con la categoría de los haces de conjuntos sobre
$\bar Y$ provistos de una acción continua de $\operatorname{Gal}(\bar k/k)$
compatible con la acción de $\operatorname{Gal}(\bar k/k)$ sobre $\bar Y$.}
\end{enumerate}

Escribamos $\bar k$ como límite directo de extensiones galoisianas finitas
$k_i$. Sean
\begin{equation*}
\varphi_i:Y_i=Y\otimes_k k_i\longrightarrow Y,
\qquad
\bar\varphi:\bar Y\longrightarrow Y
\end{equation*}
las proyecciones. Para $U$ étale sobre $Y$, supuesto afín, pongamos
$U_i=U\otimes_k k_i$. Entonces
\begin{equation*}
\bar{\mathcal F}(\bar U)
=\varinjlim_i(\varphi_i^*\mathcal F)(U_i).
\end{equation*}
Sobre $(\varphi_i^*\mathcal F)(U_i)$, el grupo
$\operatorname{Gal}(\bar k/k)$ actúa a través de
$\operatorname{Gal}(k_i/k)$, lo que demuestra (i).

El funtor de (ii) tiene por adjunto a derecha el funtor
\begin{equation*}
\mathcal G\longmapsto
(\bar\varphi_*\mathcal G)^{\operatorname{Gal}(\bar k/k)}.
\end{equation*}
Para todo $\mathcal F$ sobre $Y$ se tiene
\begin{equation*}
\bar\varphi_*\bar\varphi^*\mathcal F
=\varinjlim_i\varphi_{i*}\varphi_i^*\mathcal F.
\end{equation*}
Por tanto, para demostrar que
\begin{equation*}
\mathcal F\xrightarrow{\sim}
(\bar\varphi_*\bar\varphi^*\mathcal F)^{
  \operatorname{Gal}(\bar k/k)},
\end{equation*}
basta demostrar, para todo $i$, que
\begin{equation*}
\mathcal F\xrightarrow{\sim}
(\varphi_{i*}\varphi_i^*\mathcal F)^{
  \operatorname{Gal}(k_i/k)}.
\end{equation*}
Basta incluso demostrar que esta flecha se convierte en un isomorfismo tras
extender los escalares de $k$ a $k_i$; entonces
$\operatorname{Spec}(k_i)$ se convierte en una suma de copias de
$\operatorname{Spec}(k)$ permutadas de manera simplemente transitiva por
$\operatorname{Gal}(k_i/k)$, y la afirmación se vuelve trivial.

El funtor exacto de (ii) es, pues, plenamente fiel. Para concluir queda
comprobar que la flecha
\begin{equation*}
\bar\varphi^*(\bar\varphi_*\mathcal G)^{
  \operatorname{Gal}(\bar k/k)}\longrightarrow\mathcal G
\end{equation*}
es un epimorfismo. Sea $\bar\varphi_i:\bar Y\longrightarrow Y_i$. La
hipótesis implica que
\begin{equation*}
\varinjlim_i
\bar\varphi_i^*(\bar\varphi_{i*}\mathcal G)^{
  \operatorname{Gal}(\bar k/k_i)}
\longrightarrow\mathcal G
\end{equation*}
\pageoriginale

es un epimorfismo.

Se tiene
\begin{equation*}
\bar\varphi^*(\sisi{\varphi_*}{\bar\varphi_*}\mathcal G)^{
  \operatorname{Gal}(\bar k/k)}
=
\bar\varphi_i^*\!\left(
  \varphi_i^*\varphi_{i*}
  \left((\bar\varphi_{i*}\mathcal G)^{
    \operatorname{Gal}(\bar k/k_i)}\right)^{
    \operatorname{Gal}(k_i/k)}
\right),
\end{equation*}
y la conclusión se sigue observando, mediante el mismo argumento de descenso
galoisiano que arriba, que para $\mathcal G'$ sobre $Y_i$, provisto de una
acción de $\operatorname{Gal}(k_i/k)$ compatible con su acción sobre $Y_i$,
se tiene
\begin{equation*}
\varphi_i^*(\varphi_{i*}\mathcal G')^{
  \operatorname{Gal}(k_i/k)}
\xrightarrow{\sim}\mathcal G'.
\end{equation*}

\subsection*{\texorpdfstring{1.2. El topos $Y\times_s S$}{1.2. El topos Y por S sobre s}}

\subsubsection*{1.2.1.}
Sea $S$ un esquema local henseliano de punto cerrado
$i:s\hookrightarrow S$. El funtor
\begin{equation*}
S'\longmapsto S'\times_Ss=i^*S'
\end{equation*}
es una equivalencia de la categoría de los $S$-esquemas finitos étale $S'$ con
la categoría de los $s$-esquemas finitos étale $s'$ (espectros de
$k(s)$-álgebras étale). El funtor inverso $\mathrm{sp}^*$, del sitio étale de
$s$ al sitio étale de $S$, es un morfismo de sitios
\begin{equation*}
\mathrm{sp}:S\longrightarrow s,
\end{equation*}
el \underline{morfismo de especialización} (véase SGA~4, VIII, 7). Si $s'$ es
étale sobre $s$, se dice que $\mathrm{sp}^*(s')$ se obtiene de $S$ por
\underline{extensión del cuerpo residual}. Para un haz $F$ sobre $S$ se tiene
un isomorfismo funtorial
\begin{equation}\tag{1.2.1.1}
\mathrm{sp}_*F=i^*F.
\end{equation}
Para todo subesquema abierto $j:U\hookrightarrow S$, o para todo punto
$j:\eta\hookrightarrow S$, seguiremos escribiendo $\mathrm{sp}$ para el
morfismo compuesto $\mathrm{sp}\circ j$. Para un haz $F$ sobre $U$
(respectivamente, sobre $\eta$), se tiene
\begin{equation}\tag{1.2.1.2}
\mathrm{sp}_*F=i^*j_*F.
\end{equation}

\subsubsection*{1.2.2.}
Sea $S$ un rasgo henseliano, y sean $s,\eta,\bar s,\bar\eta,I$ como en 0.2.5.
\begin{center}
\begin{tikzcd}
s \arrow[r, hook, "i"] & S & \eta \arrow[l, hook, "j"']
\end{tikzcd}
.
\end{center}
\pageoriginale

He aquí la interpretación galoisiana de los haces sobre $s$, $\eta$ y $S$, y
del morfismo de especialización.

\begin{enumerate}[label=(\alph*)]
\item Los haces $F$ sobre $s$ (respectivamente, sobre $\eta$) se identifican,
mediante $F\mapsto F_{\bar s}$ (respectivamente,
$F\mapsto F_{\bar\eta}$), con los conjuntos provistos de una acción continua
de $\operatorname{Gal}(\bar s/s)$ (respectivamente,
$\operatorname{Gal}(\bar\eta/\eta)$) (SGA~4, VIII, 2.1, o 1.1.3).

\item Un haz $F$ sobre $S$ define haces $F_s=i^*F$ y $F_\eta=j^*F$ sobre
$s$ y $\eta$, y una flecha de adjunción $F\longrightarrow j_*j^*F$ que induce
\begin{equation*}
\varphi:F_s=i^*F\longrightarrow i^*j_*j^*F
=\mathrm{sp}_*F_\eta.
\end{equation*}
Se sabe (SGA~4, IV, 9.5.4) que el funtor
\begin{equation*}
F\longmapsto(F_s,F_\eta,\varphi)
\end{equation*}
es una equivalencia de la categoría de los haces sobre $S$ con la categoría
de las ternas $(F_s,F_\eta,\varphi)$, donde $F_s$ es un haz sobre $s$,
$F_\eta$ es un haz sobre $\eta$, y
$\varphi:F_s\longrightarrow i^*j_*F_\eta$.

Bajo la equivalencia de (a), el funtor $i^*j_*$ se identifica con el funtor
de «invariantes por $I$». Así, los haces sobre $S$ se identifican con las
ternas formadas por un conjunto $F_{\bar s}$ provisto de una acción de
$\operatorname{Gal}(\bar s/s)$, un conjunto $F_{\bar\eta}$ provisto de una
acción de $\operatorname{Gal}(\bar\eta/\eta)$, y un morfismo equivariante
$\varphi:F_{\bar s}\longrightarrow F_{\bar\eta}$.

\item Bajo (a) y (b), los morfismos $\mathrm{sp}$, $j$, $i$ y
$\mathrm{sp}=\mathrm{sp}\circ j$ se expresan como sigue:
\begin{equation*}
\begin{array}{llcl}
\mathrm{sp}:S\longrightarrow s:\;&
 \mathrm{sp}^*(F_{\bar s})&=&(F_{\bar s},F_{\bar s},\mathrm{Id}),\\[2pt]
&\mathrm{sp}_*(F_{\bar s},F_{\bar\eta},\varphi)&=&F_{\bar s},\\[3pt]
j:\eta\longrightarrow S:\;&
 j^*(F_{\bar s},F_{\bar\eta},\varphi)&=&F_{\bar\eta},\\[2pt]
&j_*(F_{\bar\eta})&=&(F_{\bar\eta}^{I},F_{\bar\eta},
 \text{inclusión}),\\[3pt]
i:s\longrightarrow S:\;&
 i^*(F_{\bar s},F_{\bar\eta},\varphi)&=&F_{\bar s},\\[2pt]
&i_*(F_{\bar s})&=&(F_{\bar s},\{0\},\text{proyección}),\\[3pt]
\mathrm{sp}=\mathrm{sp}\circ j:\eta\longrightarrow s:\;&
 \mathrm{sp}^*(F_{\bar s})&=&F_{\bar s},\\[2pt]
&\mathrm{sp}_*(F_{\bar\eta})&=&F_{\bar\eta}^{I}.
\end{array}
\end{equation*}
\end{enumerate}

\subsubsection*{1.2.3.}
Sea $S$ un rasgo henseliano como en 1.2.2, y sea $Y$ un esquema sobre $s$.
Tendremos que trabajar con el $2$-producto fibrado $Y\times_sS$ de los topoi
étale de $Y$ y $S$ sobre el topos étale de $s$. Por Giraud~[1], este
$2$-producto fibrado existe, y se comprueba fácilmente que admite la
descripción galoisiana dada a continuación. Quien prefiera evitar los
$2$-productos fibrados puede tomar esta descripción como definición.

\pageoriginale

\subsubsection*{Construcción 1.2.4.}
\underline{Los haces $F$ sobre $Y\times_sS$ se identifican con ternas
$(F_s,F_\eta,\varphi)$ en las que:}
\begin{enumerate}[label=(\alph*)]
\item $F_s$ es un haz sobre $Y$, o equivalentemente un haz $F_{\bar s}$
sobre $\bar Y=Y\otimes_s\bar s$, provisto de una acción continua de
$\operatorname{Gal}(\bar s/s)$ compatible con la acción de
$\operatorname{Gal}(\bar s/s)$ sobre $\bar Y$;

\item $F_\eta$ es un haz $F_{\bar\eta}$ sobre $\bar Y$, provisto de una
acción continua de $\operatorname{Gal}(\bar\eta/\eta)$ compatible con la
acción (mediante 1.1.1) de $\operatorname{Gal}(\bar\eta/\eta)$ sobre $\bar Y$;

\item $\varphi$ es un morfismo equivariante de $F_{\bar s}$ a
$F_{\bar\eta}$.
\end{enumerate}
Análogamente, un \underline{haz $F_\eta$ sobre $Y\times_s\eta$} es un objeto
como en 1.2.4(b). Se tiene un diagrama de topoi
\begin{equation}\tag{1.2.4.1}
\begin{tikzcd}[column sep=large, row sep=small]
Y \arrow[dd] & {} & Y \arrow[r, hook, "i"] \arrow[dd]
  & Y\times_sS \arrow[dd]
  & Y\times_s\eta \arrow[l, hook', "j"'] \arrow[dd] & {} \\
{} & {} \arrow[l, "\mathrm{sp}"'] & {} & {} & {} & {} \\
s & {} & s \arrow[r, hook, "i"] & S
  & \eta \arrow[l, hook', "j"'] & {.}
\end{tikzcd}
\end{equation}
El topos $Y\times_sS$ es la unión del subtopos abierto $Y\times_s\eta$ y su
complementario cerrado $Y$. Las fórmulas de 1.2.2(c) siguen siendo válidas
para $\mathrm{sp}:Y\times_sS\longrightarrow Y$ (también denotado
$\mathrm{pr}_1$), $j:Y\times_s\eta\hookrightarrow Y\times_sS$,
$i:Y\hookrightarrow Y\times_sS$ y
$\mathrm{sp}:Y\times_s\eta\longrightarrow Y$ (también denotado
$\mathrm{pr}_1$).

\paragraph{1.2.4.2.}
Si se toma 1.2.4 como definición de $Y\times_sS$, hay que demostrar mediante
una construcción directa que, salvo isomorfismo único, $Y\times_sS$ no
depende de la elección de una clausura separable $k(\bar\eta)$ de $k(\eta)$.
Esto se hace interpretando el $F_\eta$ de 1.2.4(b) como un funtor que, a cada
clausura separable $k(\bar\eta)$ de $k(\eta)$ que define una clausura
separable $k(\bar s)$ de $k(s)$, asocia un haz sobre
$Y\otimes_{k(s)}k(\bar s)$; el $F_\eta$ de 1.2.4(b) es la restricción de este
funtor a $\{\bar\eta\}$.

\pageoriginale

\subsubsection*{1.2.5.}
Todo punto geométrico $\bar x$ de $Y_{\bar s}$ define un punto geométrico
$(\bar x,\bar\eta)$ de $Y\times_s\eta$, cuyo funtor fibra correspondiente es
\begin{equation*}
F_\eta\longmapsto(F_{\bar\eta})_{\bar x}.
\end{equation*}
Define puntos geométricos $(\bar x,\bar\eta)$ y $(\bar x,\bar s)$ de
$Y\times_sS$ mediante
\begin{equation*}
F_{(\bar x,\bar\eta)}=(F_{\bar\eta})_{\bar x}
\qquad\text{y}\qquad
F_{(\bar x,\bar s)}=(F_{\bar s})_{\bar x}.
\end{equation*}
Se comprueba que todos los puntos de los topoi $Y\times_s\eta$ y
$Y\times_sS$ se obtienen de este modo. No utilizaremos este hecho. También se
comprueba inmediatamente que estos conjuntos de puntos son conservativos.

\subsubsection*{\texorpdfstring{${}^{*}$Observación 1.2.6.}{Observación 1.2.6.}}
La construcción de 1.2.5 define una equivalencia de la categoría
$\operatorname{Point}(Y\times_s\eta)$ (respectivamente,
$\operatorname{Point}(Y\times_sS)$) de los puntos de $Y\times_s\eta$
(respectivamente, de $Y\times_sS$) con el $2$-producto fibrado
\begin{equation*}
\operatorname{Point}(Y)\times_{\operatorname{Point}(s)}
\operatorname{Point}(\eta).
\end{equation*}
Esta equivalencia se sigue también de que $Y\times_s\eta$ (respectivamente,
$Y\times_sS$) es un $2$-producto fibrado de topoi.${}^{*}$

\subsubsection*{1.2.7.}
Por la funtorialidad de los $2$-productos fibrados, el diagrama (1.2.4.1) es
funtorial en $Y$ y $S$. En particular:
\begin{enumerate}[label=\alph*)]
\item Todo morfismo $f:Y\longrightarrow Y'$ de esquemas sobre $s$ induce un
morfismo de topoi de $Y\times_sS$ a $Y'\times_sS$. Se tiene
\begin{equation}\tag{1.2.7.1}
f^*(F_{\bar s},F_{\bar\eta},\varphi)
=(f^*F_{\bar s},f^*F_{\bar\eta},f^*(\varphi)),
\end{equation}
y, cuando $f$ es cuasicompacto,
\begin{equation}\tag{1.2.7.2}
f_*(F_{\bar s},F_{\bar\eta},\varphi)
=(f_*F_{\bar s},f_*F_{\bar\eta},f_*(\varphi)).
\end{equation}
Fórmulas análogas son válidas para
$f:Y\times_s\eta\longrightarrow Y'\times_s\eta$.

\item Todo morfismo suprayectivo de rasgos henselianos
$f:S'\longrightarrow S$ induce un morfismo de topoi de
$Y_{s'}\times_{s'}S'$ a $Y\times_sS$.

\pageoriginale

Denotando de nuevo por $f$ el morfismo
natural $Y_{s'}=Y\otimes_{k(s)}k(s')\longrightarrow Y$, se tiene
\begin{equation}\tag{1.2.7.3}
f^*(F_{\bar s},F_{\bar\eta},\varphi)
=(f^*F_{\bar s},f^*F_{\bar\eta},f^*(\varphi)).
\end{equation}
Cuando $f$ es finito, $f_*$ se expresa como una representación inducida.
\end{enumerate}

\subsubsection*{1.2.8.}
Sea $f:Y\longrightarrow Y'$ un morfismo separado, localmente de tipo finito,
entre $s$-esquemas. Se define un funtor de «imagen directa con soporte propio»
\begin{equation*}
\begin{aligned}
f_!: {}&(\text{haces de conjuntos apuntados sobre }Y\times_sS)\\
&\longrightarrow
(\text{haces de conjuntos apuntados sobre }Y'\times_sS)
\end{aligned}
\end{equation*}
mediante la fórmula
\begin{equation}\tag{1.2.8.1}
f_!(F_{\bar s},F_{\bar\eta},\varphi)
=(f_!F_{\bar s},f_!F_{\bar\eta},f_!(\varphi)).
\end{equation}
Cuando $f$ es una inmersión localmente cerrada, $f_!$ se llama funtor de
extensión por cero.

Se define un funtor análogo para $Y\times_s\eta$.

\subsubsection*{1.2.9.}
Sea $f:Y\hookrightarrow Y'$ una inmersión localmente cerrada de
$s$-esquemas. El funtor $f_!$ posee un adjunto a derecha $f^!$, «secciones con
soporte en $Y$», que satisface
\begin{equation}\tag{1.2.9.1}
f^!(F_{\bar s},F_{\bar\eta},\varphi)
=(f^!F_{\bar s},f^!F_{\bar\eta},f^!(\varphi)).
\end{equation}
Supongamos solamente que $f$ es cuasifinito. El funtor $f_!$ de los haces
abelianos sobre $Y\times_sS$ a los haces abelianos sobre $Y'\times_sS$ sigue
teniendo un adjunto a derecha $f^!$, dado por (1.2.9.1) (véase SGA~4, XVIII,
3.1.8). Si $f$ es étale, entonces $f^!=f^*$.

Se definen funtores análogos para $Y\times_s\eta$.

\pageoriginale

\subsection*{\texorpdfstring{1.3. El funtor $\Psi$}{1.3. El funtor Psi}}

\subsubsection*{1.3.1.}
Sea $S$ un rasgo henseliano como en 0.2.5, y sea
$p:X\longrightarrow S$. Nos proponemos definir un funtor exacto a izquierda
$\Psi$ de los haces sobre $X$ a los haces sobre el topos $X_s\times_sS$
construido en 1.2. En general, este funtor no es el funtor imagen directa de
un morfismo de topoi.

Definiremos también
\begin{equation*}
\Psi_\eta:(\text{haces sobre }X_\eta)
\longrightarrow(\text{haces sobre }X_s\times_s\eta).
\end{equation*}

\subsubsection*{1.3.2.}
Pongamos $\bar X=X\times_S\bar S$. Se tiene un diagrama cartesiano
\begin{equation}\tag{1.3.2.1}
\begin{tikzcd}
X_{\bar s} \arrow[r, hook, "\bar i"] \arrow[d]
  & \bar X \arrow[d]
  & X_{\bar\eta} \arrow[l, hook', "\bar j"'] \arrow[d] \\
X_s \arrow[r, hook] & X & X_\eta \arrow[l, hook']
\end{tikzcd}
\end{equation}
sobre el diagrama
\begin{center}
\begin{tikzcd}
\bar s \arrow[r, hook] \arrow[d]
  & \bar S \arrow[d]
  & \bar\eta \arrow[l, hook'] \arrow[d] \\
s \arrow[r, hook] & S & \eta \arrow[l, hook']
\end{tikzcd}
.
\end{center}
Sea $F$ un haz sobre $X_\eta$, con imagen inversa $F_{\bar\eta}$ sobre
$X_{\bar\eta}$. Pongamos
\begin{equation}\tag{1.3.2.2}
\Psi_\eta(F)=\bar i^*\bar j_*F_{\bar\eta}.
\end{equation}
Por transporte de estructura, este haz está provisto de una acción continua
de $\operatorname{Gal}(\bar\eta/\eta)$ compatible con la acción de
$\operatorname{Gal}(\bar\eta/\eta)$ sobre $X_{\bar s}$. El funtor
\begin{equation}\tag{1.3.2.3}
\Psi_\eta:(X_{\eta,\mathrm{et}})^{\sim}
\longrightarrow(X_s\times_s\eta)^{\sim}
\end{equation}
es exacto a izquierda.

\subsubsection*{1.3.3.}
Sea $F$ un haz sobre $X$, y sean $F_\eta$ y $F_s$ sus restricciones a
$X_\eta$ y $X_s$. Pongamos
\pageoriginale

\begin{equation}\tag{1.3.3.1}
(\Psi(F))_\eta=\Psi_\eta(F_\eta)
\end{equation}
y
\begin{equation}\tag{1.3.3.2}
(\Psi(F))_s=F_s.
\end{equation}
Sea $\bar F$ la imagen inversa de $F$ sobre $\bar X$, y sea $\varphi'$ el
morfismo de adjunción
\begin{equation*}
\varphi':\bar F\longrightarrow\bar j_*\bar j^*\bar F.
\end{equation*}
Este morfismo induce
\begin{equation*}
\varphi:F_{\bar s}=\bar i^*\bar F
\longrightarrow
\bar i^*\bar j_*\bar j^*\bar F
=(\Psi_\eta(F_\eta))_{\bar\eta}.
\end{equation*}
La terna
\begin{equation*}
\Psi(F)=((\Psi(F))_s,(\Psi(F))_\eta,\varphi)
\end{equation*}
es un haz sobre $X_s\times_sS$. El funtor
\begin{equation}\tag{1.3.3.3}
\Psi:(X_{\mathrm{et}})^{\sim}
\longrightarrow(X_s\times_sS)^{\sim}
\end{equation}
es exacto a izquierda.

\subsubsection*{1.3.4.}
Cuando no haya riesgo de confusión, escribiremos
\begin{equation*}
\Psi_\eta(F)=(\Psi(F))_\eta,
\qquad
\Psi_{\bar\eta}(F)=(\Psi(F))_{\bar\eta},
\qquad\text{y}\qquad
\Psi_{\bar\eta}(F)=(\Psi_\eta(F))_{\bar\eta}.
\end{equation*}

\subsubsection*{1.3.5.}
En lo que resta de esta sección enumeramos propiedades de funtorialidad de
$\Psi$ (véanse 1.2.7--1.2.9). Serán útiles sobre todo en su forma derivada
(2.1).

Sea $f:X\longrightarrow X'$ un morfismo de $S$-esquemas, que da un diagrama
\begin{equation}\tag{1.3.5.1}
\begin{tikzcd}
X_{\bar s} \arrow[r, hook, "\bar i"] \arrow[d, "f"]
  & \bar X \arrow[d, "f"]
  & X_{\bar\eta} \arrow[l, hook', "\bar j"'] \arrow[d, "f"] \\
X'_{\bar s} \arrow[r, hook, "\bar i"]
  & \bar X' & X'_{\bar\eta} \arrow[l, hook', "\bar j"']
\end{tikzcd}
\end{equation}
(la barra de su $\bar X$ central superior es visible en el impreso). La
ambigüedad de la notación de las flechas simplifica las fórmulas
y no causa confusión: solo una interpretación da sentido a las composiciones
que siguen. Asimismo, escribimos $\Psi$ para el funtor (1.3.3.3) tanto para
$X$ como para $X'$.

Los morfismos de funtores que se describen a continuación para $\Psi$ poseen
variantes para $\Psi_\eta$, que por lo general omitimos enunciar
explícitamente.

\pageoriginale

\subsubsection*{1.3.6.}
El morfismo de cambio de base
$\bar i^*f_*\longrightarrow f_*\bar i^*$ induce un morfismo de funtores
\begin{equation}\tag{1.3.6.1}
\Psi f_*\longrightarrow f_*\Psi,
\end{equation}
que es un isomorfismo cuando $f$ es propio (SGA~4, XII, 5.1(i)).

El caso esencial es $X'=S$. Se obtiene un morfismo
\begin{equation}\tag{1.3.6.2}
f_*\longrightarrow f_*\Psi
\end{equation}
cuyo componente esencial es
\begin{equation}\tag{1.3.6.3}
\Gamma(X_{\bar\eta},F)
\longrightarrow
\Gamma(X_{\bar s},\Psi_{\bar\eta}(F));
\end{equation}
este morfismo es un isomorfismo cuando $f$ es propio.

\subsubsection*{1.3.7.}
El morfismo de cambio de base
$f^*\bar j_*\longrightarrow\bar j_*f^*$ induce un morfismo de funtores
\begin{equation}\tag{1.3.7.1}
f^*\Psi\longrightarrow\Psi f^*.
\end{equation}

\subsubsection*{1.3.8.}
Supongamos que $f$ es cuasifinito. Los funtores $f_!$ conmutan con el cambio
de base; el morfismo $f_!\bar j_*\longrightarrow\bar j_*f_!$ induce un
morfismo de funtores
\begin{equation}\tag{1.3.8.1}
f_!\Psi\longrightarrow\Psi f_!.
\end{equation}
Cuando $f$ es finito, este es el inverso de (1.3.6.1).

\subsubsection*{1.3.9.}
La situación es dual para $f^!$ (suponemos que $f$ es cuasifinito al trabajar
con haces abelianos, o que es una inmersión al trabajar con haces de conjuntos
apuntados). Este funtor conmuta con las imágenes directas, y el morfismo
$\bar i^*f^!\longrightarrow f^!\bar i^*$ induce un morfismo de funtores
\begin{equation}\tag{1.3.9.1}
\Psi f^!\longrightarrow f^!\Psi.
\end{equation}
Cuando $f$ es étale, este es el inverso de (1.3.7.1).

\pageoriginale

\subsubsection*{1.3.10.}
Por último, un cambio de base $S'\xrightarrow{f}S$ da un morfismo de cambio
de base
\begin{equation}\tag{1.3.10.1}
f^*\Psi\longrightarrow\Psi f^*.
\end{equation}
\subsection*{1.4. La variación}

\subsubsection*{1.4.1.}
Sea $S$ un rasgo henseliano como en 0.2.5, y sea $Y$ un esquema sobre $s$.
Sea $\Lambda$ un anillo. Más generalmente, podría tomarse como $\Lambda$ un
haz de anillos sobre $Y$. Para todo objeto $K$ de la categoría derivada
$D(Y\times_sS,\Lambda)$ de la categoría de los haces de $\Lambda$-módulos
sobre $Y\times_sS$, nos proponemos definir un morfismo de
\underline{variación}.

\subsubsection*{1.4.2.}
Un complejo de haces de $\Lambda$-módulos $K$ sobre
$\sisi{X}{Y}\times_sS$ se identifica con una terna
$(K_s,K_\eta,\varphi)$ en la que
\begin{enumerate}[label=\alph*)]
\item $K_s$ (respectivamente, $K_\eta$) es un complejo de haces
$K_{\bar s}$ (respectivamente, $K_{\bar\eta}$) sobre $\bar Y$, provisto de
una acción continua de $\operatorname{Gal}(\bar s/s)$ (respectivamente,
$\operatorname{Gal}(\bar\eta/\eta)$);

\item $\varphi$ es un morfismo equivariante
$\varphi:K_{\bar s}\longrightarrow K_{\bar\eta}$.
\end{enumerate}
Aquí y en lo que sigue mantenemos fijo el esquema $Y$ de 1.4.1; la frase
inicial impresa de 1.4.2 escribe inadvertidamente $X\times_sS$.

Un complejo $K$ de este tipo es siempre homotópico a un complejo
$K'=(K'_s,K'_\eta,\varphi')$ tal que $\varphi'$ es inyectivo y la sucesión
exacta
\begin{equation}\tag{1.4.2.1}
0\longrightarrow K'_{\bar s}
\longrightarrow K'_{\bar\eta}
\longrightarrow\operatorname{coker}(\varphi')
\longrightarrow0
\end{equation}
es escindida grado por grado. Basta tomar $K'_s=K_s$ y tomar como $K'_\eta$
la suma directa de $K_\eta$ y el cono de $\mathrm{sp}^*K_s$. Pongamos
\begin{equation*}
\Phi(K)=\operatorname{coker}(\varphi').
\end{equation*}
La sucesión exacta (1.4.2.1) define un triángulo distinguido
\begin{equation}\tag{1.4.2.2}
\begin{tikzcd}[column sep=large, row sep=large]
{} & \Phi(K) \arrow[dl] & {} \\
\mathrm{sp}^*K_s \arrow[rr] & {} & K_\eta \arrow[ul]
\end{tikzcd}
\end{equation}
en la categoría $K(Y\times_s\eta,\Lambda)$ de los complejos de haces de
$\Lambda$-módulos sobre $Y\times_s\eta$, salvo homotopía. Este triángulo solo
depende de $K\in\operatorname{Ob}K(Y\times_sS,\Lambda)$.

\pageoriginale

La construcción pasa a la categoría derivada y asocia a todo objeto $K$ de
$D(Y\times_sS,\Lambda)$ un triángulo distinguido de
$D(Y\times_s\eta,\Lambda)$.

El triángulo (1.4.2.2) es útil sobre todo por la sucesión exacta larga de
hipercohomología que induce:
\begin{equation*}
\cdots\longrightarrow H^i(\bar Y,K_{\bar s})
\longrightarrow H^i(\bar Y,K_{\bar\eta})
\longrightarrow H^i(\bar Y,\Phi(K)_{\bar\eta})
\longrightarrow\cdots.
\end{equation*}

\subsubsection*{1.4.3.}
El grupo de inercia $I$ actúa trivialmente sobre $K'_{\bar s}$. Por tanto,
para todo $\sigma\in I$, el endomorfismo $\sigma-1$ de $K'_{\bar\eta}$
(notación (1.4.2.1)) se factoriza a través de un morfismo
\begin{equation*}
\operatorname{Var}(\sigma):
\operatorname{coker}(\varphi')\longrightarrow K'_{\bar\eta}.
\end{equation*}
Esta construcción pasa a la categoría derivada: para
$K\in\operatorname{Ob}D(Y\times_sS,\Lambda)$ y $\sigma\in I$, define la
\underline{variación}
\begin{equation}\tag{1.4.3.1}
\operatorname{Var}(\sigma):
\Phi(K)_{\bar\eta}\longrightarrow K_{\bar\eta}.
\end{equation}
Denotando por $q$ la flecha de $K_\eta$ a $\Phi(K)$, se tiene
\begin{align}
\sigma&=1+\operatorname{Var}(\sigma)q
&&\text{sobre }K_{\bar\eta}, \tag{1.4.3.2}\\
\sigma&=1+q\operatorname{Var}(\sigma)
&&\text{sobre }\Phi(K)_{\bar\eta}, \tag{1.4.3.3}
\end{align}
y
\begin{equation}\tag{1.4.3.4}
\operatorname{Var}(\sigma\tau)
=\operatorname{Var}(\sigma)q\operatorname{Var}(\tau)
 +\operatorname{Var}(\sigma)+\operatorname{Var}(\tau).
\end{equation}

\section*{2. Ciclos evanescentes}

\subsection*{\texorpdfstring{2.1. El funtor $R\Psi$}{2.1. El funtor R Psi}}

\subsubsection*{2.1.1.}
Sea $S$ un rasgo henseliano como en 0.2.5, y sea $\Lambda$ un anillo de
torsión cuya torsión es prima a la característica residual $p$ de $S$
(0.2.7). Más generalmente, podría tomarse como $\Lambda$ un haz de torsión
cuya torsión fuera prima a las características residuales; el caso más
importante es $\Lambda=\mathbb Z/\ell^n$, con $\ell$ un primo invertible en
$\mathcal O_S$.

\pageoriginale

Sea $X$ un esquema sobre $S$. El funtor $\Psi$ (respectivamente,
$\Psi_\eta$) lleva haces de $\Lambda$-módulos sobre $X$ (respectivamente,
sobre $X_\eta$) a haces de $\Lambda$-módulos sobre $X_s\times_sS$
(respectivamente, sobre $X_s\times_s\eta$). Sean $R\Psi$ y $R\Psi_\eta$ los
funtores derivados
\begin{equation*}
R\Psi:D^+(X,\Lambda)\longrightarrow D^+(X_s\times_sS,\Lambda),
\end{equation*}
\begin{equation*}
R\Psi_\eta:D^+(X_\eta,\Lambda)
\longrightarrow D^+(X_s\times_s\eta,\Lambda),
\end{equation*}
y sea $R\Phi$ el funtor compuesto $\Phi\circ R\Psi$:
\begin{equation*}
R\Phi:D^+(X,\Lambda)
\longrightarrow D^+(X_s\times_s\eta,\Lambda).
\end{equation*}

\subsubsection*{2.1.2.}
Consideremos el diagrama conmutativo (1.3.2.1)
\begin{equation}\tag{1.3.2.1}
\begin{tikzcd}
X_{\bar s} \arrow[r, hook, "\bar i"] \arrow[d]
  & \bar X \arrow[d]
  & X_{\bar\eta} \arrow[l, hook', "\bar j"'] \arrow[d] \\
X_s \arrow[r, hook, "i"] & X & X_\eta \arrow[l, hook', "j"']
\end{tikzcd}
\end{equation}
Si $F$ es un haz inyectivo de módulos sobre $X$, su restricción al subesquema
abierto $X_\eta$ sigue siendo inyectiva. Si $F_\eta$ es inyectivo sobre
$X_\eta$, su imagen inversa $F_{\bar\eta}$ es acíclica. Por tanto, las formas
derivadas de (1.3.3.1) y (1.3.3.2) son
\begin{equation}\tag{2.1.2.1}
R\Psi(K)_s=i^*K
\qquad(K\in\operatorname{Ob}D^+(X,\Lambda)),
\end{equation}
\begin{equation}\tag{2.1.2.2}
R\Psi(K)_\eta=R\Psi_\eta(K_\eta)
\qquad(K\in\operatorname{Ob}D^+(X,\Lambda)),
\end{equation}
y
\begin{equation}\tag{2.1.2.3}
R\Psi_\eta(K)_{\bar\eta}
=\bar i^*R\bar j_*K_{\bar\eta}
\qquad(K\in\operatorname{Ob}D^+(X_\eta,\Lambda)).
\end{equation}
Esto permite escribir, sin conflicto,
\begin{equation*}
R\Psi_\eta(K)=(R\Psi(K))_\eta,
\qquad
R\Psi_{\bar\eta}(K)=(R\Psi(K))_{\bar\eta},
\qquad\text{y}\qquad
R\Psi_{\bar\eta}(K)=(R\Psi_\eta(K))_{\bar\eta}.
\end{equation*}
\pageoriginale

Con esta notación, el triángulo distinguido (1.4.2.2) se convierte en
\begin{equation}\tag{2.1.2.4}
\begin{tikzcd}[row sep=large, column sep=huge]
{} & R\Phi(K) \arrow[dl] & {} \\
\mathrm{sp}^*i^*K \arrow[rr] & {}
  & R\Psi_\eta(K_{\sisi{s}{\eta}}) \arrow[ul]
\end{tikzcd}
.
\end{equation}
El término inferior derecho impreso contiene $K_s$; la identidad definitoria
(2.1.2.2) exige $K_\eta$, como se muestra en la rama corregida.

Para $\sigma\in I$ se tiene un morfismo de \underline{variación}
\begin{equation}\tag{2.1.2.5}
\operatorname{Var}(\sigma):R\Phi(K)_{\bar\eta}
\longrightarrow R\Psi_{\bar\eta}(K).
\end{equation}
Los haces $R^i\Phi(\Lambda)$, o más generalmente $R^i\Phi(K)$, son los
\underline{haces de ciclos evanescentes}.

\subsubsection*{2.1.3.}
Sea $(\bar x,\bar\eta)$ el punto geométrico de $X_s\times_s\eta$ definido
por un punto geométrico $\bar x$ de $X_{\bar s}$ (1.2.5). Calculamos la fibra
de $R\Psi_\eta(K_\eta)$ en $(\bar x,\bar\eta)$ mediante (2.1.2.3) (véase I,
2.3). El henselizado estricto $X_{(\bar x)}$ de $X$ en $\bar x$ es un esquema
sobre el henselizado estricto $S_{\mathrm{nr}}$ de $S$, con punto geométrico
genérico $\eta_{\mathrm{nr}}$.

\subsubsection*{Proposición 2.1.4.}
\underline{Se tiene}
\begin{equation*}
(R\Psi(K_\eta))_{(\bar x,\bar\eta)}
=R\Gamma\!\left(
X_{(\bar x)}\times_{\eta_{\mathrm{nr}}}\bar\eta,K
\right);
\end{equation*}
\underline{en particular,}
\begin{equation*}
R^i\Psi(K_\eta)_{(\bar x,\bar\eta)}
=\mathbb H^i\!\left(
X_{(\bar x)}\times_{\eta_{\mathrm{nr}}}\bar\eta,K
\right).
\end{equation*}

Apliquemos SGA~4, XV, 2.1 (el teorema de aciclicidad local de los morfismos
lisos). Se obtiene la reformulación siguiente.

\subsubsection*{Reformulación 2.1.5.}
\underline{Si $f$ es liso y $F$ es un haz localmente constante, entonces}
\begin{equation*}
R\Phi(F)=0.
\end{equation*}
Más generalmente, si $K\in\operatorname{Ob}D^+(X,\Lambda)$, entonces
$R\Phi(K)=0$ allí donde $f$ es liso y los haces $\mathcal H^i(K)$ son
localmente constantes. Véase también 2.1.11.

\subsubsection*{2.1.6.}
Sea $f:Y\longrightarrow Y'$ un morfismo de esquemas sobre $s$. Los funtores
$f^*$, $f_*$, $f^!$ y $f_!$ (1.2.7--1.2.9) poseen análogos naturales en la
categoría derivada.
\begin{enumerate}[label=\alph*)]
\item El funtor $f^*$ es exacto y, por tanto, se deriva trivialmente.

\pageoriginale

\item El funtor $f_*$ posee un funtor derivado a derecha $Rf_*$. Cuando $f$
es cuasicompacto, se tiene
\begin{equation}\tag{2.1.6.1}
(Rf_*K)_{\bar s}=Rf_{\bar s*}(K_{\bar s})
\end{equation}
y
\begin{equation}\tag{2.1.6.2}
(Rf_*K)_{\bar\eta}=Rf_{\bar s*}(K_{\bar\eta}).
\end{equation}

\item Cuando $f$ es cuasifinito, el funtor $f_!$ es exacto y, por tanto, se
deriva trivialmente. En general, el funtor derivado de $f_!$ es patológico, y
se procede como en SGA~4, XVII. Supongamos que $f$ es separado y de tipo
finito y que $Y'$ es noetheriano. Existe entonces [2] una factorización
$f=\widetilde f\,j$, donde $j$ es una inmersión abierta y $\widetilde f$ es
propio, y se define
\begin{equation*}
Rf_!=R\widetilde f_*\circ j_!.
\end{equation*}
Los argumentos de SGA~4, XVII muestran que $Rf_!$ es independiente de la
elección de la factorización $\widetilde f\,j$. Los análogos de (2.1.6.1) y
(2.1.6.2) son válidos.

\item Utilizaremos $Rf^!$ únicamente para un morfismo cuasifinito, en cuyo
caso es el funtor derivado de $f^!$. También aquí son válidos los análogos de
(2.1.6.1) y (2.1.6.2).
\end{enumerate}

\subsubsection*{2.1.7.}
Revisemos ahora las propiedades de funtorialidad de $R\Psi$, es decir,
derivemos 1.3.5--1.3.10. Sea $f:X\longrightarrow X'$ un morfismo de
$S$-esquemas, y conservemos la notación de 1.3.5.

El morfismo (1.3.6.1) define
\begin{equation}\tag{2.1.7.1}
R\Psi\,Rf_*\longrightarrow Rf_*\,R\Psi.
\end{equation}
Por el teorema de cambio de base propio (SGA~4, XII, 5.1),
\underline{este morfismo es un isomorfismo si $f$ es propio.}

El morfismo (1.3.7.1) define
\begin{equation}\tag{2.1.7.2}
f^*R\Psi\longrightarrow R\Psi f^*.
\end{equation}
Por el teorema de cambio de base liso (SGA~4, XVI, 1.2),
\underline{este morfismo es un isomorfismo si $f$ es liso.}

\pageoriginale

El isomorfismo (2.1.7.1) para $f$ propio y el morfismo
(\sisi{2.1.8.1}{1.3.8.1}) definen un morfismo
\begin{equation}\tag{2.1.7.3}
Rf_!R\Psi\longrightarrow R\Psi Rf_!,
\end{equation}
que, cuando $f$ es propio, es el inverso de (2.1.7.1). La impresión remite
por error al caso particular posterior (2.1.8.1).

Cuando $f$ es cuasifinito, (1.3.9.1) se deriva en
\begin{equation}\tag{2.1.7.4}
R\Psi Rf^!\longrightarrow Rf^!R\Psi,
\end{equation}
que, cuando $f$ es étale, es el inverso de
(\sisi{2.1.7.1}{2.1.7.2}). La referencia impresa a (2.1.7.1) es incompatible
con los dos funtores implicados.

Sea $f:S'\longrightarrow S$ un cambio de rasgo, y pongamos
$X'=X\times_SS'$:
\begin{center}
\begin{tikzcd}
X' \arrow[r, "f"] \arrow[d, "p"'] & X \arrow[d, "p"] \\
S' \arrow[r, "f"] & S
\end{tikzcd}
\end{center}
(para la ambigüedad de la notación, véase 1.3.5). Supongamos que $f$ es
suprayectivo. La flecha de cambio de base (1.3.10.1) se deriva en
\begin{equation}\tag{2.1.7.5}
f^*R\Psi\longrightarrow R\Psi f^*.
\end{equation}
Las mismas afirmaciones son válidas para $R\Psi_\eta$.

\subsubsection*{2.1.8.}
El caso esencial de (2.1.7.1) y (2.1.7.3) es $X'=S$. Para
$K\in\operatorname{Ob}D^+(X_{\bar\eta},\Lambda)$ se obtienen entonces
morfismos
\begin{equation}\tag{2.1.8.1}
Rf_*\longrightarrow Rf_*R\Psi
\end{equation}
y
\begin{equation}\tag{2.1.8.2}
Rf_!R\Psi\longrightarrow Rf_!,
\end{equation}
de donde se deducen
\begin{equation}\tag{2.1.8.3}
\mathbb H^i(X_{\bar\eta},K)
\longrightarrow
\mathbb H^i(X_{\bar s},R\Psi_{\bar\eta}(K))
\end{equation}
y
\begin{equation}\tag{2.1.8.4}
\mathbb H_c^i(X_{\bar s},R\Psi_{\bar\eta}(K))
\longrightarrow
\mathbb H_c^i(X_{\bar\eta},K),
\end{equation}
así como la conmutatividad de los diagramas

\pageoriginale

\begin{equation}\tag{2.1.8.5}
\begin{tikzcd}[column sep=large]
(R^if_*K)_{\bar s} \arrow[r, "\mathrm{sp}"] \arrow[d]
  & \mathbb H^i(X_{\bar\eta},K) \arrow[d] \\
\mathbb H^i(X_{\bar s},K) \arrow[r]
  & \mathbb H^i(X_{\bar s},R\Psi_{\bar\eta}(K))
\end{tikzcd}
\end{equation}
y
\begin{equation}\tag{2.1.8.6}
\begin{tikzcd}[column sep=large]
\mathbb H_c^i(X_{\bar s},K) \arrow[rr, "\mathrm{sp}"] \arrow[dr]
  & {} & \mathbb H_c^i(X_{\bar\eta},K) \\
{} & \mathbb H_c^i(X_{\bar s},R\Psi_{\bar\eta}(K)) \arrow[ur] & {}
\end{tikzcd}
.
\end{equation}

Para $K\in\operatorname{Ob}D^+(X_\eta,\Lambda)$ y $\sigma\in I$, los
diagramas siguientes son conmutativos:
\begin{equation}\tag{2.1.8.7}
\begin{tikzcd}[column sep=small, cells={nodes={font=\small}}]
\mathbb H^i(X_{\bar\eta},K)
  \arrow[rr, dash, "\sigma-1"] \arrow[d]
  & {} & \mathbb H^i(X_{\bar\eta},K) \arrow[d] \\
\mathbb H^i(X_{\bar s},R\Psi_{\bar\eta}(K)) \arrow[r]
  & \mathbb H^i(X_{\bar s},R\Phi_{\bar\eta}(K))
      \arrow[r, "\operatorname{Var}(\sigma)"]
  & \mathbb H^i(X_{\bar s},R\Psi_{\bar\eta}(K))
\end{tikzcd}
\end{equation}
y
\begin{equation}\tag{2.1.8.8}
\begin{tikzcd}[column sep=small, cells={nodes={font=\small}}]
\mathbb H_c^i(X_{\bar\eta},K)
  \arrow[rr, "\sigma-1"]
  & {} & \mathbb H_c^i(X_{\bar\eta},K) \\
\mathbb H_c^i(\sisi{X_s}{X_{\bar s}},R\Psi_{\bar\eta}(K)) \arrow[r] \arrow[u]
  & \mathbb H_c^i\bigl(X_{\bar s},R\Phi_{\bar\eta}(K)\sisi{\bigr.}{\bigr)}
      \arrow[r, "\operatorname{Var}(\sigma)"]
  & \mathbb H_c^i(X_{\bar s},R\Psi_{\bar\eta}(K)) \arrow[u]
\end{tikzcd}
.
\end{equation}
En (2.1.8.8), la rama corregida restituye la barra omitida en el término
inferior izquierdo impreso $X_s$ y el paréntesis de cierre que falta en
el término inferior central. En (2.1.8.7), en cambio, el conector superior
está impreso realmente como un trazo sin punta de flecha, mientras que
(2.1.8.8) tiene una flecha; se conserva esta distinción visible en la fuente.

Cuando $f$ es propio, (2.1.8.1) y (2.1.8.2), así como (2.1.8.3) y (2.1.8.4),
son isomorfismos mutuamente inversos. El triángulo distinguido (2.1.2.4)
proporciona una sucesión exacta larga
\begin{equation}\tag{2.1.8.9}
\cdots\longrightarrow\mathbb H^i(X_{\bar s},K)
\longrightarrow\mathbb H^i(X_{\bar\eta},K)
\longrightarrow\mathbb H^i(X_{\bar s},R\Phi_{\bar\eta}(K))
\longrightarrow\cdots.
\end{equation}
Así, el complejo evanescente $R\Phi(K)$ expresa la diferencia entre las
cohomologías de las fibras geométricas $X_{\bar s}$ y $X_{\bar\eta}$, con
coeficientes en la imagen inversa de $K$. Por (2.1.8.7), la variación
$\operatorname{Var}(\sigma)$ determina la acción del grupo de inercia $I$
sobre $\mathbb H^i(X_{\bar\eta},K)$. Estos hechos sustentan todas las
aplicaciones del método de los ciclos evanescentes.

\pageoriginale

La hipótesis de propiedad impuesta arriba a $f$ puede reemplazarse por
condiciones de regularidad en el infinito sobre $f$ y $K$. Por ejemplo, se
tiene el resultado siguiente (para
$K\in\operatorname{Ob}D(X_{\bar\eta},\Lambda)$).

\subsubsection*{Proposición 2.1.9.}
\underline{Supongamos que $f$ admite una factorización $f=f_1k$}
\begin{center}
\begin{tikzcd}[column sep=huge]
X \arrow[r, hook, "k"] & X_1 \arrow[r, "f_1"] & S,
\end{tikzcd}
\end{center}
\underline{donde $f_1$ es propio y $k$ identifica $X$ con el complementario
en $X_1$ de un divisor de cruces normales relativo $D$} (por tanto, por
hipótesis, $f_1$ es liso en un entorno de $D$). \underline{Si, en un entorno
$U$ de $D$, los haces $\mathcal H^i(K)$ son localmente constantes sobre
$U\setminus(D\cup X_0)$ y moderadamente ramificados a lo largo del divisor de
cruces normales $D\cup X_0$, entonces (2.1.8.1) y (2.1.8.2) son
isomorfismos.}

La sucesión espectral de hipercohomología de $K$ muestra que basta tratar el
caso en que $K$ consiste en un haz $F_\eta$ situado en grado cero. Tras
reemplazar $S$ por su normalización en una extensión finita (lo cual está
permitido, como se ve directamente a partir de las definiciones) y aplicar el
lema de Abhyankar, podemos suponer que $F_\eta$ es la restricción a $X_\eta$
de un haz $F$ sobre $X$, localmente constante sobre $U\cap X$ y moderadamente
ramificado a lo largo de $D$. Con la notación de 2.1.2, sea $\overline F$ la
imagen inversa de $F$ sobre $\overline X$, y escribamos
$\bar f:\overline X\longrightarrow\overline S$.

Para demostrar que (2.1.8.1) es un isomorfismo basta demostrar que
\[
R\bar f_*\bigl(R\bar j_*\bar j^*\overline F\bigr)
\]
conmuta con el cambio de base por $\bar s\longrightarrow\overline S$. En el
triángulo distinguido
\[
\longrightarrow\overline F\longrightarrow
R\bar j_*\bar j^*\overline F\longrightarrow\Delta\longrightarrow,
\]
los haces $\mathcal H^i(\Delta)$ tienen soporte en $X_{\bar s}$; para
$\Delta$, la afirmación de cambio de base en cuestión es tautológica. Queda,
pues, demostrar que $R\bar f_*(\overline F)$ conmuta con el cambio de base por
$\bar s\longrightarrow\overline S$, o bien, teniendo en cuenta el teorema de
cambio de base para el morfismo propio $f_1$, que $Rk_*(\overline F)$ conmuta
con el cambio de base $X_{1\bar s}\longrightarrow X_1$. Este es el Lema
2.1.10 siguiente.

\pageoriginale

Para demostrar que (2.1.8.2) es un isomorfismo basta demostrar, para el haz
$F$, que la extensión por cero $k_!$ conmuta con el funtor $R\bar j_*$; es
decir, en un entorno de $D$,
\[
R\Phi(k_!F)=0\qquad\text{sobre }X_{1\bar s}.
\]
Este es el Lema 2.1.11 siguiente.

\subsubsection*{Lema 2.1.10.}
\underline{Sea $f:X\longrightarrow S$ liso, sea $D\subset X$ un divisor de
cruces normales relativo, pongamos
$U=X-D\mathrel{\mathop{\hookrightarrow}^{j}}X$, y sea $F$ un haz de
$\Lambda$-módulos sobre $U$, localmente constante y moderadamente ramificado
a lo largo de $D$. La formación de $Rj_*F$ conmuta con todo cambio de base
$T\longrightarrow S$.}

La cuestión es local sobre $X$ para la topología étale. Podemos suponer, por
tanto, que existe un morfismo $g:X'\longrightarrow X$, donde $X'$ es la
normalización de $X$ en un recubrimiento finito étale suprayectivo
$g':U'\longrightarrow U$, moderadamente ramificado a lo largo de $D$, tal que
$g'^*F$ es constante sobre $U'$. Por el lema de Abhyankar, podemos suponer
además que $X'$ es liso sobre $S$ y que $U'$ es el complementario en $X'$ de
un divisor de cruces normales relativo $D'$.

El haz $F$ se inyecta en $g'_*g'^*F$, y el conúcleo vuelve a satisfacer las
hipótesis del lema. Iterando esta construcción (el mismo $g$ sigue sirviendo)
se obtiene una resolución de $F$ por haces de la forma $g'_*G$, donde $G$ es
constante y $g'$ es como arriba. Basta, pues, demostrar la afirmación para tal
haz, o equivalentemente para $(X'/S,D',G)$ con $G$ constante sobre $U'$. Este
es un caso particular de SGA~5, II, 4.4.

\subsubsection*{Lema 2.1.11.}
\underline{Sea $S$ un rasgo henseliano, sea $f:X\longrightarrow S$ liso, sea
$D\subset X$ un divisor de cruces normales relativo, pongamos
$U=X-D\mathrel{\mathop{\hookrightarrow}^{j}}X$, y sea $F$ un haz de
$\Lambda$-módulos sobre $U$, localmente constante y moderadamente ramificado
a lo largo de $D$. Entonces}
\[
\underline{R\Phi(j_!F)=0.}
\]

\pageoriginale

El dévissage utilizado acaba de reducirnos al caso en que $F$ es constante,
de valor $G$. Podemos suponer que $D$ es la unión de divisores lisos
$(D_i)_{i\in I}$. Para $J\subset I$, sea $D_J$ la intersección de los $D_i$
con $i\in J$. El haz $F$ posee una resolución $F^\bullet$ cuyos términos son
sumas de haces $G_J$, constantes sobre $D_J$ y prolongados por cero a $X$.
Basta, por tanto, demostrar el lema para uno de tales haces, y esto se sigue
de (2.1.5), aplicado a $D_J/S$.

Conjeturo que (2.1.7.5) es un isomorfismo cuando $S$ y $S'$ son excelentes.
He aquí un resultado parcial (véase también 2.4.2).

\subsubsection*{Proposición 2.1.12.}
\underline{El morfismo (2.1.7.5) es un isomorfismo en los casos siguientes:}

\medskip
\noindent(a) $S$ es excelente y $S'$ es su completado;

\noindent(b) $S$ y $S'$ tienen característica igual a cero.

En el caso (b), nos reducimos a la situación en que $S$ y $S'$ son
estrictamente locales y un uniformizante de $S$ sigue siendo un uniformizante
de $S'$ (se extienden los escalares sin cambiar
$\bar S'\longrightarrow\bar S$). En ambos casos (a) y (b), denotando con una
barra la normalización en una clausura separable del cuerpo de fracciones, se
tiene entonces
\[
\bar S'\simeq S'\times_S\bar S.
\]
Volviendo a las definiciones, basta demostrar que $S'\longrightarrow S$ es
universalmente localmente acíclico (aplíquese SGA~4, XVI, 1.1). Esto se sigue
de I, 0.5.1 (basado en la versión de Artin de la desingularización de Néron) y
SGA~4, XV, 2.1 (el teorema de aciclicidad local para morfismos lisos).

\subsubsection*{2.1.13.}
Supongamos que $f:X\longrightarrow S$ es de tipo finito. La dimensión
cohomológica (SGA~4, XVII, 1.2.9) de $R\Psi$ es entonces a lo sumo
$\dim(X_\eta)$, como se deduce de 2.1.4 y SGA~4, XIV, 3.2 (véase I, 4.2(i)).
Por tanto, $R\Psi$ se extiende a toda la categoría derivada,
\[
R\Psi:D(X,\Lambda)\longrightarrow D(X_s\times_sS,\Lambda),
\]
y lleva $D^-$ en $D^-$. Sea $L\in\operatorname{Ob}D^-(\Lambda)$ un complejo
acotado superiormente de $\Lambda$-módulos a derecha, y sea
$K\in\operatorname{Ob}D^-(X,\Lambda)$. Se tiene entonces un isomorfismo

\pageoriginale

\begin{equation}\tag{2.1.13.1}
L\overset{\mathbf L}{\otimes}_\Lambda R\Psi(K)
\xrightarrow{\ \sim\ }
R\Psi\bigl(L\overset{\mathbf L}{\otimes}_\Lambda K\bigr).
\end{equation}
Para comprobar (2.1.13.1), se reemplaza primero $L$ por un complejo de
$\Lambda$-módulos libres. Como $R\Psi$ es triangulado y de dimensión
cohomológica finita, basta comprobar la afirmación después de reemplazar $L$
por uno de sus términos. El paso al colímite filtrante nos reduce al caso en
que $L$ es un módulo libre finito, y luego al caso tautológico $L=\Lambda$.

Se deduce de (2.1.13.1), mediante el argumento habitual (véase SGA~4, XVII,
5.2.11), que si $K\in\operatorname{Ob}D(X_\eta,\Lambda)$ tiene dimensión de
Tor finita a lo sumo $d$, donde $d\in\mathbb Z$ (SGA~4, XVII, 4.1.9),
entonces $R\Psi_\eta(K)$ también tiene dimensión de Tor a lo sumo $d$.

\subsection*{2.2. Una compatibilidad (morfismos traza)}

\subsubsection*{2.2.1.}
Sea $S$ un rasgo henseliano, y sea
$u:X_s\longrightarrow Y_s$ un morfismo separado, cuasifinito, plano y de
presentación finita entre $s$-esquemas. Para todo haz abeliano $F$ sobre
$Y_s$ se tiene un morfismo traza
\[
\operatorname{Tr}_u:u_!u^*F\longrightarrow F
\]
(SGA~4, XVII, 6.2.3). Si $F$ es un haz abeliano sobre $Y_s\times_sS$ o sobre
$Y_s\times_s\eta$, se obtiene un morfismo análogo
$\operatorname{Tr}_u:u_!u^*F\longrightarrow F$.

Sea ahora $u:X\longrightarrow Y$ un morfismo separado, cuasifinito, plano y
de presentación finita entre $S$-esquemas. Seguimos escribiendo $u$ para el
morfismo inducido
\[
X_s\times_sS\longrightarrow Y_s\times_sS,
\]
y utilizamos $\Psi$ para los funtores de 1.3.3, tanto para $X$ como para $Y$.

\subsubsection*{Proposición 2.2.2.}
\underline{Para todo haz abeliano $F$ sobre $Y$, el diagrama siguiente es
conmutativo:}
\begin{center}
\begin{tikzcd}[column sep=large, row sep=large]
u_!u^*\Psi F
  \arrow[r, "{(1.3.7.1)}"] \arrow[d, "\operatorname{Tr}_u"']
  & u_!\Psi u^*F \arrow[r, "{(1.3.8.1)}"]
  & \Psi u_!u^*F \arrow[d, "\operatorname{Tr}_u"] \\
\Psi F \arrow[rr, equal] & {} & \Psi F.
\end{tikzcd}
\end{center}
\pageoriginale

La verificación se deja al lector. Hay una variante para $\Psi_\eta$ y para
$F$ sobre $Y_\eta$.

\subsubsection*{2.2.3.}
Sea $\Lambda$ como en \sisi{2.2.1}{2.1.1}, y sea
$u:X\longrightarrow Y$ un morfismo cuasifinito y plano entre $S$-esquemas
separados, de tipo finito y planos:
\begin{center}
\begin{tikzcd}[column sep=large, row sep=large]
X \arrow[rr, "u"] \arrow[dr, "f"'] & {} & Y \arrow[dl, "g"] \\
{} & S & {}
\end{tikzcd}
\end{center}
La impresión remite aquí a 2.2.1, pero el anillo de coeficientes $\Lambda$ se
fija en 2.1.1. Seguimos escribiendo $f$ y $g$ para los morfismos
\[
X_s\times_sS\longrightarrow S,
\qquad
Y_s\times_sS\longrightarrow S,
\]
y suponemos que $\dim Y_\eta\leq n$. Para todo
$K\in\operatorname{Ob}D(S,\Lambda)$ existe un morfismo traza (SGA~4, XVIII,
2.9)
\[
\operatorname{Tr}_f:Rf_!f^*K(n)[2n]\longrightarrow K.
\]

\subsubsection*{Corolario 2.2.4.}
\underline{El diagrama siguiente es conmutativo:}
\begin{center}
\resizebox{\textwidth}{!}{%
\begin{tikzcd}[ampersand replacement=\&, column sep=large, row sep=large]
Rg_!u_!u^*R\Psi\,g^*K(n)[2n]
  \arrow[r, "{(1.3.7.1)}"] \arrow[d, "\operatorname{Tr}_u"']
  \& Rf_!R\Psi\,f^*K(n)[2n]
      \arrow[r, "{(2.1.10.1)}"]
  \& Rf_!f^*K(n)[2n] \arrow[d, "\operatorname{Tr}_f"] \\
Rg_!R\Psi\,g^*K(n)[2n]
  \arrow[r, "{(2.1.10.1)}"]
  \& Rg_!g^*K(n)[2n] \arrow[r, "\operatorname{Tr}_g"]
  \& K
\end{tikzcd}}
\end{center}
\underline{Si $u$ es étale, (1.3.7.1) es un isomorfismo.}

Omitiendo los argumentos, el diagrama (2.2.4) se desarrolla como sigue:
\begin{center}
\resizebox{\textwidth}{!}{%
\begin{tikzcd}[ampersand replacement=\&, column sep=large, row sep=huge]
Rg_!u_!u^*R\Psi\,g^*
  \arrow[r] \arrow[d, "\operatorname{Tr}_u"']
  \arrow[rd, phantom, "{(1)}" description]
  \& Rg_!R\Psi\,u_!u^*g^*
      \arrow[r] \arrow[d, "\operatorname{Tr}_u"']
      \arrow[rd, phantom, "{(2)}" description]
  \& Rg_!u_!u^*g^*
      \arrow[r, "\sim"] \arrow[d, "\operatorname{Tr}_u"']
      \arrow[rd, phantom, "{(3)}" description]
  \& Rf_!f^* \arrow[d, "\operatorname{Tr}_f"] \\
Rg_!R\Psi\,g^* \arrow[r, equal]
  \& Rg_!R\Psi\,g^* \arrow[r]
  \& Rg_!g^* \arrow[r, "\operatorname{Tr}_g"]
  \& K.
\end{tikzcd}}
\end{center}
\pageoriginale

La conmutatividad de (1) se sigue de 2.2.2; la de (2) es la funtorialidad de
(2.1.10.2), y la de (3) es la transitividad de la traza (SGA~4, XVIII, 2.9,
Variante~3).

\subsubsection*{2.2.4.1.}
Hay un resultado análogo para $\Psi_\eta$ y
$K\in\operatorname{Ob}D(\eta,\Lambda)$.

\subsubsection*{2.2.5.}
Sea $\Phi$ un subconjunto de $Y_s$ propio sobre $s$, y denotemos por
$Rg_\Phi$ el funtor «imagen directa con soporte en $\Phi$». Sea
$\operatorname{Tr}_g$ el compuesto
\begin{equation}\tag{2.2.5.1}
\begin{split}
\operatorname{Tr}_g:
\sisi{Rf_\Phi}{Rg_\Phi} R\Psi g^*K(n)[2n]
&\longrightarrow Rg_!R\Psi g^*K(n)[2n] \\
&\longrightarrow Rg_!g^*K(n)[2n]
\longrightarrow K,
\end{split}
\end{equation}
y análogamente para $\Psi_\eta$ y $K\in D(\eta,\Lambda)$. La impresión
contiene $Rf_\Phi$ como primer funtor de (2.2.5.1); la definición
inmediatamente anterior y el resto del compuesto exigen $Rg_\Phi$, que se usa
en la rama corregida.

Si $\Phi$ es un conjunto finito de puntos cerrados, entonces
\[
Rg_\Phi\circ u_!=Rf_{u^{-1}\Phi},
\]
y por tanto existe un morfismo traza
\[
\operatorname{Tr}_u:
Rf_{u^{-1}\Phi}u^*=Rg_\Phi u_!u^*\longrightarrow Rg_\Phi.
\]
Por 2.2.4, el diagrama siguiente es conmutativo:
\begin{center}
\resizebox{\textwidth}{!}{%
\begin{tikzcd}[ampersand replacement=\&, column sep=large, row sep=large]
Rf_{u^{-1}\Phi}u^*R\Psi\,g^*K(n)[2n]
  \arrow[r] \arrow[d, "\operatorname{Tr}_u"']
  \& Rf_{u^{-1}\Phi}R\Psi\,f^*K(n)[2n]
      \arrow[r, "\operatorname{Tr}_f"]
  \& K \arrow[d, equal] \\
Rg_\Phi R\Psi\,g^*K(n)[2n]
  \arrow[rr, "\operatorname{Tr}_g"']
  \& {} \& K.
\end{tikzcd}}
\end{center}

Supongamos que $u$ es étale, sea $x\in X_s$, y supongamos que $x$ e
$y=u(x)$ tienen el mismo cuerpo residual. Entonces
\[
Rf_{\{x\}}\xrightarrow{\ \sim\ }Rg_{\{\sisi{x}{y}\}}.
\]
La impresión contiene $\{x\}$ también a la derecha, pero el soporte para $g$
está en $Y_s$ y es $\{y\}$, como confirma el diagrama siguiente. Este
isomorfismo, dado por $\operatorname{Tr}_u$, expresa la localidad étale de la
cohomología con soporte en $x$. El diagrama
\begin{equation}\tag{2.2.5.2}
\begin{tikzcd}[column sep=huge, row sep=large]
Rf_{\{x\}}R\Psi\,f^*K(n)[2n]
  \arrow[r, "\operatorname{Tr}_f"] \arrow[d, equal]
  & K \arrow[d, equal] \\
Rg_{\{y\}}R\Psi\,g^*K(n)[2n]
  \arrow[r, "\operatorname{Tr}_g"]
  & K
\end{tikzcd}
\end{equation}

\pageoriginale

es conmutativo.

\subsubsection*{Escolio 2.2.6.}
La compatibilidad (\sisi{2.3.5.2}{2.2.5.2}) dice que el morfismo traza
(\sisi{2.3.5.1}{2.2.5.1}), para $\Phi=\{\sisi{x}{y}\}$, es un invariante
local: solo depende del henselizado del $S$-esquema $X$ en $x$. Lo mismo vale
para $R\Psi_\eta$. Las referencias impresas a (2.3.5.2) y (2.3.5.1) son las
inmediatamente anteriores (2.2.5.2) y (2.2.5.1); asimismo,
$\Phi\subset Y_s$ hace de $\{y\}$, y no del $\{x\}$ impreso, el soporte
pertinente.

\subsection*{2.3. Teorema de finitud (en característica igual a cero)}

\subsubsection*{Teorema 2.3.1.}
\underline{Sea $S$ un rasgo henseliano como en 0.2.5, sea $\Lambda$ un anillo
noetheriano de torsión y sea $f:X\longrightarrow S$ un morfismo de tipo
finito. Supongamos que $S$ tiene característica igual a cero. Entonces, para
todo haz constructible de $\Lambda$-módulos $F$ sobre $X_\eta$, los haces
$R^i\Psi_{\bar\eta}(F)$ sobre $X_{\bar s}$ son constructibles.}

Comenzamos repitiendo el principio de la demostración de SGA~4, XIV, 1.10.
Apliquemos SGA~4, IX, 2.14(ii). Existe una resolución $F^\bullet$ de $F$
\begin{center}
\begin{tikzcd}[column sep=large, row sep=large]
{} & F \arrow[d] & {} & {} \\
0 \arrow[r] & F^0 \arrow[r] & F^1 \arrow[r] & \cdots
\end{tikzcd}
\end{center}
por haces de la forma $\bigoplus_i p'_{i*}G_i$, donde
$p'_i:Y_{i\eta}\longrightarrow X_\eta$ es finito y $G_i$ es un haz
constructible constante. Podemos suponer reducido $Y_{i\eta}$; sea
$p_i:Y_i\longrightarrow X$ la normalización de $X$ en $Y_{i\eta}$. Mediante
la sucesión espectral de la resolución $F^\bullet$ y aplicando (2.1.7.1) en el
caso elemental del morfismo finito $p_i$, se ve que basta demostrar el teorema
para cada par $(Y_i,G_i)$.

Demostremos a continuación el teorema cuando $F$ es el haz constante definido
por un $\Lambda$-módulo $M$. Recordemos que $\Lambda$ es una
$\mathbb Z/n$-álgebra para cierto $n$. El módulo $M$ puede, por tanto,
descomponerse por dévissage en módulos $M_i$ sobre
$\Lambda/\ell_i\Lambda$, donde cada $\ell_i\mid n$ es un primo adecuado.
Esto nos reduce al caso en que $\Lambda$ es una $\mathbb Z/\ell$-álgebra, con
$\ell$ primo. La fórmula (2.1.13.1) da entonces
\[
M\otimes_{\mathbb Z/\ell}
R^i\Psi_{\bar\eta}(\mathbb Z/\ell)
\xrightarrow{\ \sim\ }
R^i\Psi_{\bar\eta}(M).
\]
\pageoriginale

Queda por tratar el haz constante $F=\mathbb Z/\ell$.

Elijamos una sucesión decreciente de subconjuntos cerrados
\[
X_\eta=X'_0\supset\cdots\supset X'_n
\]
tal que $(X'_i-X'_{i+1})_{\mathrm{red}}$ sea liso sobre $\eta$, y sea
$X_i=\overline{X'_i}$ la clausura en $X$. El haz constante
$\mathbb Z/\ell$ se descompone por dévissage en los haces constantes
$\mathbb Z/\ell$ sobre los estratos $X'_i-X'_{i+1}$, prolongados por cero.
Así:

\medskip
\noindent\underline{Reducción 1.}
Basta demostrar el Teorema 2.3.1 para un haz $j_!\mathbb Z/\ell$, donde
$j:U\hookrightarrow X_\eta$ es la inclusión de un abierto liso.

Por resolución de singularidades existe un morfismo
$g:X_1\longrightarrow X$ tal que

\medskip
\noindent(a) $g$ es propio e induce un isomorfismo
$g^{-1}(U)\xrightarrow{\sim}U$;

\noindent(b) $X_1$ es regular y $X_1-g^{-1}(U)$ es un divisor de cruces
normales de $X_1$, unión de divisores regulares $D_\alpha$.

Sea $j_1:g^{-1}(U)\hookrightarrow X_1$. Entonces
\[
Rg_*j_{1!}\mathbb Z/\ell=j_!\mathbb Z/\ell.
\]
Aplicando (2.1.7.1) y el teorema de finitud SGA~4, XIV, 1.1, se ve que basta
demostrar el Teorema 2.3.1 para $(X_1,j_{1!}\mathbb Z/\ell)$.

Existe una resolución
\[
0\longrightarrow j_{1!}\mathbb Z/\ell
\longrightarrow\mathbb Z/\ell
\longrightarrow\bigoplus_\alpha(\mathbb Z/\ell)_{D_\alpha}
\longrightarrow
\bigoplus_{\alpha,\beta}
(\mathbb Z/\ell)_{D_\alpha\cap D_\beta}
\longrightarrow\cdots,
\]
y por tanto:

\medskip
\noindent\underline{Reducción 2.}
Basta demostrar el Teorema 2.3.1 para el haz constante $\mathbb Z/\ell$
sobre un $S$-esquema regular $X$ tal que $X_s$ sea un divisor de cruces
normales en $X$.

En este último caso se utiliza el cálculo explícito I, 3.3, que a su vez usa
el teorema de pureza SGA~4, XIX, 3.2 y la hipótesis de característica cero.

\subsection*{2.4. Singularidades aisladas}

\subsubsection*{2.4.1.}
En toda esta sección suponemos, por simplicidad, que $S$ es un rasgo
estrictamente henseliano. Sea $\Lambda$ un anillo noetheriano de torsión cuya
torsión es prima a la característica residual de $S$, sea
$f:X\longrightarrow S$ un morfismo de tipo finito, y sea $F$ un haz
constructible de $\Lambda$-módulos sobre $X$. Un punto $x\in X$ se llama

\pageoriginale

punto liso de $(X,F)$ si $f$ es liso en $x$ y $F$ es localmente constante en
un entorno de $x$. Sea $\Sigma$ el conjunto de los puntos donde $(X,F)$ no
es liso, y sea $\Sigma_s$ su fibra especial. Sabemos que $R\Phi(F)=0$ fuera
de $\Sigma_s$.

\subsubsection*{Teorema 2.4.2.}
\underline{En todo punto aislado $x$ de $\Sigma_s$, los módulos
$R^i\Phi_{\bar\eta}(F)_x$ son $\Lambda$-módulos de tipo finito e invariantes
por cambio de rasgo (mediante 2.1.7.1).}

En característica igual a cero, estas afirmaciones se demostraron en 2.1.12
y 2.3.1 sin suponer que $x$ fuera aislado.

Sea $x$ un punto aislado de $\Sigma_s$. El teorema es local en un entorno de
$x$, por lo que podemos suponer que $f$ es propio. Se tiene entonces la
sucesión exacta larga (2.1.8.9)
\[
\cdots\longrightarrow H^i(X_s,F)
\longrightarrow H^i(X_{\bar\eta},F)
\longrightarrow
\mathbb H^i\bigl(X_s,R\Phi_{\bar\eta}(F)\bigr)
\longrightarrow\cdots .
\]
En esta sucesión, $H^i(X_s,F)$ y $H^i(X_{\bar\eta},F)$ son invariantes por
cambio de rasgo y son $\Lambda$-módulos de tipo finito (SGA~4, XIV, 1.1). Los
módulos $\mathbb H^i(X_s,R\Phi_{\bar\eta}(F))$ tienen, por tanto, las mismas
propiedades. Recordemos el lema siguiente.

\subsubsection*{Lema 2.4.3.}
\underline{Sea $(T,\Lambda)$ un topos anillado, sea $Z$ un subtopos cerrado
de $T$, y sea $i:Z\hookrightarrow T$ la inclusión. Entonces el funtor}
\[
i_*:D(Z,\Lambda)\longrightarrow D(T,\Lambda)
\]
\underline{es una equivalencia de $D(Z,\Lambda)$ sobre la subcategoría plena
de $D(T,\Lambda)$ formada por los complejos $K$ para los que
$\mathcal H^i(K)$ tiene soporte en $Z$.}

En efecto, para tal complejo $K$, el morfismo
$K\longrightarrow i_*i^*K$ es un cuasi-isomorfismo.

En consecuencia, podemos considerar $R\Phi_{\bar\eta}(F)$ como un objeto de
\[
D^b(\Sigma_s,\Lambda)
=D^b(\Sigma_s\setminus\{x\},\Lambda)\oplus D^b(\{x\},\Lambda).
\]
Pasando a un sumando directo, se obtiene que
\[
\mathbb H^i\bigl(\{x\},R\Phi_{\bar\eta}(F)\bigr)
=R^i\Phi_{\bar\eta}(F)_x
\]
posee las propiedades enunciadas en 2.4.2. Esto termina la demostración.

\subsubsection*{Observación 2.4.4.}
La misma demostración muestra que, si $K$ es constructible en
$D^+(X,\Lambda)$ y $x$ está aislado en el soporte de
$R\Phi_{\bar\eta}(K)$, entonces los módulos de cohomología del tallo
$R\Phi_{\bar\eta}(K)_x$ son de tipo finito sobre $\Lambda$. Si, después de
un cambio de rasgo, $x$ sigue siendo un punto aislado del soporte, estos
módulos son invariantes por dicho cambio de rasgo.

\pageoriginale

\subsubsection*{2.4.5.}
Sea $x\in X_s$ un punto aislado donde $(X,F)$ no es liso. Más generalmente,
sea $K\in\operatorname{Ob}D^b(X,\Lambda)$, y sea $x$ un punto aislado del
soporte de $R\Phi(K)$. Sea $i:\{x\}\hookrightarrow X_s$ la inclusión. El
complejo $i_*i^*R\Phi(K)$, que es isomorfo a $R\Phi(K)$ en un entorno de $x$,
es un sumando directo de $R\Phi(\sisi{F}{K})$ por 2.4.3. Así, para
$\sigma\in I$, el morfismo de variación induce
\[
\operatorname{Var}(\sigma):
i_*i^*R\Phi_{\bar\eta}(K)\longrightarrow R\Psi_{\bar\eta}(K),
\]
es decir,
\begin{equation}\tag{2.4.5.1}
\operatorname{Var}(\sigma):
\bigl(R\Phi_{\bar\eta}(K)\bigr)_x
\longrightarrow
R\Gamma_{\{x\}}\bigl(R\Psi_{\bar\eta}(K)\bigr).
\end{equation}
El $R\Phi(F)$ impreso en la frase sobre el sumando directo es $R\Phi(K)$ en
esta formulación general.

\subsubsection*{2.4.6.}
Supongamos que el soporte $\Sigma$ de $R\Phi(K)$ está formado por un número
finito de puntos aislados. Para todo $\sigma\in I$, el endomorfismo $\sigma-1$
de $R\Psi_{\bar\eta}(K)$ es entonces el compuesto representado en el diagrama
siguiente:
\begin{equation}\tag{2.4.6.1}
\resizebox{\textwidth}{!}{%
\begin{tikzcd}[ampersand replacement=\&, column sep=huge, row sep=huge]
R\Psi_{\bar\eta}(K)
  \arrow[r] \arrow[d, "\sigma-1"']
  \& R\Phi_{\bar\eta}(K)
     =\displaystyle\bigoplus_{x\in\Sigma}
       i_{x*}\bigl(R\Phi_{\bar\eta}(K)\bigr)_x
     \arrow[d, "{\bigoplus i_{x*}\operatorname{Var}(\sigma)}"'] \\
R\Psi_{\bar\eta}(K)
  \& \displaystyle\bigoplus_{x\in\Sigma}
       i_{x*}R\Gamma_{\{x\}}R\Psi_{\bar\eta}(K)
     \arrow[l]
\end{tikzcd}}
\end{equation}

Si $f$ es propio, al aplicar $H^i$ a (2.4.6.1) se obtiene que el endomorfismo
$\sigma-1$ de $\mathbb H^i(X_{\bar\eta},K)$ es el compuesto
\begin{equation}\tag{2.4.6.2}
\resizebox{\textwidth}{!}{%
\begin{tikzcd}[ampersand replacement=\&, column sep=large, row sep=huge]
\mathbb H^i(X_{\bar\eta},K)
  \arrow[r, "\sim"] \arrow[d, "\sigma-1"']
  \& \mathbb H^i\bigl(X_s,R\Psi_{\bar\eta}(K)\bigr)
     \arrow[r]
  \& \mathbb H^i\bigl(X_s,R\Phi_{\bar\eta}(K)\bigr)
     =\displaystyle\bigoplus_{x\in\Sigma}R^i\Phi_{\bar\eta}(K)_x
     \arrow[d, "{\bigoplus\operatorname{Var}(\sigma)}"'] \\
\mathbb H^i(X_{\bar\eta},K)
  \& \mathbb H^i\bigl(X_s,R\Psi_{\bar\eta}(K)\bigr)
     \arrow[l, "\sim"']
  \& \displaystyle\bigoplus_{x\in\Sigma}
     \mathbb H^i_{\{x\}}\bigl(X_s,R\Psi_{\bar\eta}(K)\bigr)
     \arrow[l]
\end{tikzcd}}
\end{equation}
Esto describe la acción de $I$ sobre $\mathbb H^i(X_{\bar\eta},K)$; sus
ingredientes esenciales son los morfismos de variación (2.4.5.1), que son
locales sobre $X$. También se tiene la sucesión exacta larga

\pageoriginale

\begin{equation}\tag{2.4.6.3}
\longrightarrow \mathbb H^i(X_s,K)
\longrightarrow \mathbb H^i(X_{\bar\eta},K)
\longrightarrow \bigoplus_{x\in\Sigma}R^i\Phi_{\bar\eta}(K)_x
\longrightarrow\cdots .
\end{equation}

El mismo argumento se aplica cuando $f$ no es necesariamente propio, siempre
que $(X,K)$ satisfaga 2.1.9. El diagrama (2.4.6.5) siguiente define un
morfismo
\begin{equation}\tag{2.4.6.4}
\operatorname{Var}(\sigma):
R\Gamma(X_{\bar\eta},K)\longrightarrow R\Gamma_c(X_{\bar\eta},K)
\end{equation}
tal que el endomorfismo $\sigma-1$ de $R\Gamma(X_{\bar\eta},K)$
(respectivamente, de $R\Gamma_c(X_{\bar\eta},K)$) es el compuesto
$\alpha\operatorname{Var}(\sigma)$ (respectivamente,
$\operatorname{Var}(\sigma)\alpha$), donde
\[
\alpha:R\Gamma_c(X_{\bar\eta},K)\longrightarrow
R\Gamma(X_{\bar\eta},K).
\]
\begin{equation}\tag{2.4.6.5}
\resizebox{\textwidth}{!}{%
\begin{tikzcd}[ampersand replacement=\&, column sep=huge, row sep=huge]
R\Gamma(X_{\bar\eta},K)
  \arrow[r, "\sim"] \arrow[d, "\operatorname{Var}(\sigma)"']
  \& R\Gamma\bigl(X_s,R\Psi_{\bar\eta}(K)\bigr)
     \arrow[r]
  \& \displaystyle\bigoplus_{x\in\Sigma}
     \bigl(R\Phi_{\bar\eta}(K)\bigr)_x
     \arrow[d, "\bigoplus"']
  \& {} \\
R\Gamma_c(X_{\bar\eta},K)
  \arrow[r, "\sim"']
  \& R\Gamma_c\bigl(X_s,R\Psi_{\bar\eta}(K)\bigr)
  \& \displaystyle\bigoplus_{x\in\Sigma}
     R\Gamma_{\{x\}}\bigl(X_s,R\Psi_{\bar\eta}(K)\bigr)
     \arrow[l]
  \& .
\end{tikzcd}}
\end{equation}

En este caso se tienen asimismo las sucesiones exactas largas
\begin{equation}\tag{2.4.6.6}
\resizebox{\textwidth}{!}{%
\begin{tikzcd}[ampersand replacement=\&, column sep=large, row sep=large]
\cdots \arrow[r]
  \& \mathbb H^i_c(X_s,K) \arrow[r] \arrow[d]
  \& \mathbb H^i_c(X_{\bar\eta},K) \arrow[r] \arrow[d]
  \& \displaystyle\bigoplus_{x\in\Sigma}
     \bigl(R^i\Phi_{\bar\eta}(K)\bigr)_x
     \arrow[r] \arrow[d, equal]
  \& \cdots \\
\cdots \arrow[r]
  \& \mathbb H^i(X_s,K) \arrow[r]
  \& \mathbb H^i(X_{\bar\eta},K) \arrow[r]
  \& \displaystyle\bigoplus_{x\in\Sigma}
     R^i\Phi_{\bar\eta}(K)_x \arrow[r]
  \& \cdots
\end{tikzcd}}
\end{equation}

La referencia impresa (\sisi{2.4.4.1}{2.4.5.1}) en la discusión que sigue a
(2.4.6.2) es el morfismo de variación (2.4.5.1).

\subsubsection*{2.4.6. (segunda aparición en la fuente)}
Sea $X_{s(x)}$ el henselizado de $X_s$ en $x$, y pongamos
\[
X^*_{s(x)}=X_{s(x)}\setminus\{x\}.
\]
Las sucesiones exactas largas de hipercohomología de
$(X_{s(x)},\{x\})$, con coeficientes en $K_s$,
$R\Psi_{\bar\eta}(K)$ y $R\Phi_{\bar\eta}(K)$, dan el siguiente diagrama
conmutativo de sucesiones exactas largas:

\pageoriginale

\begin{equation}\tag{2.4.6.1}
\resizebox{\textwidth}{!}{%
\begin{tikzcd}[ampersand replacement=\&, column sep=large, row sep=large]
{}
  \& 0 \arrow[r] \arrow[d]
  \& R^{n-1}\Phi_{\bar\eta}(K)_x \arrow[r] \arrow[d]
  \& R^{n-1}\Phi_{\bar\eta}(K)_x \arrow[r] \arrow[d]
  \& 0 \arrow[d]
  \& {} \\
{} \arrow[r]
  \& \mathbb H^{n-1}(X^*_{s(x)},K)
     \arrow[r, "j"] \arrow[d, equal]
  \& \mathbb H^n_{\{x\}}(X_s,K) \arrow[r] \arrow[d]
  \& \underline{H}^{n}(K)_x \arrow[r] \arrow[d]
  \& \mathbb H^n(X^*_{s(x)},K)
     \arrow[r] \arrow[d, equal]
  \& {} \\
{} \arrow[r]
  \& \mathbb H^{n-1}(X^*_{s(x)},K)
     \arrow[r, "j", "{\mbox{\textcircled{1}}}"'] \arrow[d]
  \& \mathbb H^n_{\{x\}}\bigl(X_s,R\Psi_{\bar\eta}(K)\bigr)
     \arrow[r] \arrow[d]
  \& R^n\Psi_{\bar\eta}(K)_x \arrow[r] \arrow[d]
  \& \mathbb H^n(X_{s(x)},K) \arrow[r] \arrow[d]
  \& {} \\
{}
  \& 0 \arrow[r]
  \& R^n\Phi_{\bar\eta}(K)_x \arrow[r, equal]
  \& R^n\Phi_{\bar\eta}(K)_x \arrow[r]
  \& 0
  \& {}
\end{tikzcd}}
\end{equation}
La imagen de la flecha \textcircled{1} está contenida en los invariantes por
la inercia.

La fuente imprime tanto este párrafo como el precedente con el número 2.4.6,
y ambos diagramas principales con el número (2.4.6.1). Se conservan aquí los
números repetidos; esta segunda aparición es la citada en el Exposé XV,
2.2.6.

\renewcommand{\refname}{Bibliografía}
\begin{thebibliography}{9}
\bibitem{GiraudClassifyingTopos}
J.~Giraud, «Classifying topos», pp.~43--56 en
\emph{Toposes, Algebraic Geometry and Logic},
Lecture Notes in Mathematics 274, Springer-Verlag, 1971.

\bibitem{NagataImbedding}
M.~Nagata, «A generalization of the imbedding problem»,
\emph{J. Math. Kyoto} \textbf{3}, n.º~1 (1963), pp.~89--102.
\end{thebibliography}
